\section{Implications for Individuals, Organizations, and Society}
\label{sec:implications}

A measured rate is only useful if it changes what somebody does, and the measurement here is
that no revenue shortfall is established on either tax object. This section says what
follows for the people a labor-to-capital shift reaches, for the organizations that would have
to act, and for the public interest. We mark what the evidence supports and what it does not,
because the gap we measure is large enough to invite recommendations it cannot carry.

\paragraph{For the people a shift reaches.} The household result is the one with the most direct
bearing on individuals, and it is unwelcome. Capacity to service existing debt falls for about
\result{HHBottomQuartileAffected} percent of bottom-wealth-quartile households against about
\result{HHTopOnePctAffected} percent of the top one percent, monotone in wealth. The reason is
not complicated: the households that owe the debt are not the households that own the equity, and
contracts do not move when factor shares do. A payment fixed in nominal terms becomes harder to
meet exactly when the income behind it falls.

The remedy that suggests itself, giving households a share of the capital, does not work in the
form usually proposed. Broadening ownership through retirement vehicles moves the boundary by
exactly zero unless the holdings can be spent, because ownership without access cannot service a
payment. For the individual at the point of decision this is the practical statement: a
retirement account is not a buffer, and policy that treats broadened ownership as though it were
one will not reach the person whose payment is due this month. Access, not ownership, is the
operative variable.

\paragraph{For managers and organizations.} The conclusion for anyone planning around a capital
tax instrument is that the instrument is weaker than its headline rate suggests, and weak for
structural reasons that will not be fixed by a rate change. An organization modeling its own
exposure to a tax rise on capital income should model the owner layer explicitly rather than
applying a statutory rate, because the owner layer adds only \result{ShareholderLayerPct} percent
to the rate on rents at central values.

There is a timing point as well, and it runs against the optimistic reading. Under full expensing
the deduction precedes the income, so while investment is growing, current deductions exceed the
tax on income from a smaller earlier capital stock. Revenue actually collected during an
investment boom is below what the assembled rate implies, not above it. Any plan that counts on
the boom paying for the transition has the sign wrong.

\paragraph{Strategic implications.} The ownership finding reframes a debate that has been
conducted largely in terms of enforcement. If the binding constraint were profit shifting, the
instruments would be information exchange, transfer pricing rules and minimum taxes, and the
problem would be one of administrative capacity. Our ranking says the constraint is third on that
list. The instruments that actually move the rate are the ones that touch deferral, step-up at
death and the boundary of the taxable shareholder base, which are domestic, legislative and
politically harder, and which have had a fraction of the attention.

The second strategic point is about hedging, and it cuts against the state's own position. The
rate links the two directions of risk without trading one against the other: revenue at risk in a
bust is linear in the capital tax rate, so raising the rate enough to close the displacement-side
condition scales bust-state exposure by the same multiple. A government cannot tax its way out of
one direction without increasing what it loses in the other. That is an argument for instruments
that are not rates at all, and we say so without claiming to know which.

\paragraph{For society and equity.} The same concentration of ownership runs through every result
in this paper, and its distributional reading is blunt. The capital tax reaches little because
most corporate claims sit where the income tax does not. Households cannot be made whole by
ownership because most corporate equity is not theirs and what is theirs cannot be spent. A bust
transmits weakly to households because the assets that fall are held by the households least
likely to change their spending. One fact about who owns American corporate capital explains a
revenue problem, an incidence problem and a transmission result.

That is a statement about the structure of claims, not about anybody's deserts, and we draw one
narrow conclusion from it. A shift of income from labor to capital is not distributionally
neutral and cannot be made so by adjusting a rate after the fact, because the rate operates on a
base that most households are not in. Whether to act on that, and how, is a question for a
legislature. What the measurement supplies is the constraint any answer has to satisfy, and the
knowledge that the obvious answers, a higher rate and broader nominal ownership, do not satisfy
it.
